Fraud screening systems must rank suspicious transactions quickly and tell the analyst why a transaction was flagged. Graph neural networks (GNNs) use the relations between transactions but usually need separate post-hoc explainers, whose output is costly to compute and not guaranteed to reflect the model. We propose RTXGNN, a self-explaining temporal GNN. Learned feature, edge and node masks gate its message passing and are returned as the explanation of every score. A training objective makes the top-$k$ explanation sufficient for the prediction and its complement uninformative. A shift-invariant temporal encoding (HRAPE) supplies time gaps, recency and transaction velocity without encoding absolute time. On the Elliptic Bitcoin dataset with a temporal split, and against 16 baselines over 10 seeds, RTXGNN reaches a test F1 of \rtxFone{} and an average precision of \rtxAP{}. This is on par with the strongest GNN baselines (best: \bestgnnName, F1 \bestgnnFone) but not significantly better after correction for multiple comparisons. Tree ensembles on the same features are equally strong (random forest \rfFone, XGBoost \xgbFone). At the chosen threshold the model still misses \rtxMissPct\% of illicit transactions, and all models fail after a regime change in the test period; retraining with a one-step label delay recovers part of this loss. On the same trained model, RTXGNN's explanations are more faithful than saliency, Integrated Gradients and GNNExplainer, and they come at no extra inference cost. On a synthetic benchmark with known laundering rings they identify the planted transactions (edge AUC \edgeAucdefaultrtxgnnseal{}). Scoring and explaining one transaction end to end, including neighbourhood sampling, takes \rtxLatBoneP{} ms at the 99th percentile on a single CPU core. The explanations can support analyst review and audit trails. They do not by themselves establish regulatory compliance.
